% Source: https://artofproblemsolving.com/community/c3h1352656_degree_of_a_monomial
% Stable identity: aops-topic-1352656-q01
Add the exponents of all the variables (a variable written without an exponent has exponent $1$, and a constant has no variables).
\begin{enumerate}
\item[(a)] $3+3=\boxed{6}$. (It is not $3$: the degree is the total of the exponents, not the largest one. For example, $x^3y^3=(xy)^3$ behaves like a sixth-degree quantity, since replacing $x,y$ by $tx,ty$ multiplies it by $t^6$.)
\item[(b)] $2+1+5=\boxed{8}$; the coefficient $-4$ does not affect the degree.
\item[(c)] A nonzero constant has degree $\boxed{0}$.
\end{enumerate}
